% ==============================================================================
% 封面页（LaTeX + TikZ 忠实复刻 1985 年原版《形形色色的曲线》封面）
% ==============================================================================
\clearpage
\thispagestyle{empty}
\null

% 原书封面色彩精确取样定义
\definecolor{origcoverbg}{RGB}{102, 127, 116}
\definecolor{origcovergreen}{RGB}{156, 224, 166}
\definecolor{origcoverdark}{RGB}{48, 65, 57}
\definecolor{origivory}{RGB}{253, 254, 244}
\definecolor{origivoryshade}{RGB}{241, 237, 226}
\definecolor{origcurve}{RGB}{68, 102, 83}

\begin{tikzpicture}[remember picture, overlay]
  % 1. 纯净复古松石灰绿底色（原扫描件左侧暗痕为纸张压褶，非色块）
  \fill[origcoverbg] (current page.south west) rectangle (current page.north east);

  \begin{scope}[shift={(current page.south west)}, x=1cm, y=1cm]
    \clip (0, 0) rectangle (13.0, 18.4);

    % 2. 巨型亮草绿鹦鹉螺对数螺线（单条主螺线略粗 + 单向 C 形气室隔壁）
    \coordinate (naut) at (10.10, 16.05);
    \begin{scope}[shift={(naut)}]
      % 单条外延主螺线：r(t) = 0.0685 * exp(0.1685 * t)，顺时针角度 t
      \draw[line width=2.85pt, origcovergreen, line cap=round, smooth, domain=1.0:28.9, samples=440]
        plot ({0.0685 * exp(0.1685*\x) * cos(\x r)}, {-0.0685 * exp(0.1685*\x) * sin(\x r)});

      % 均匀密集的 C 形弧形气室隔壁：起于内圈单螺线、顺时针微弓后平滑扫掠接至外圈单螺线
      \foreach \k in {0,...,49} {
        \pgfmathsetmacro{\th}{10.2 + \k * 0.345}
        \pgfmathsetmacro{\rOut}{0.0685 * exp(0.1685*\th)}
        \pgfmathsetmacro{\rIn}{0.0685 * exp(0.1685*(\th - 2*pi))}
        \pgfmathsetmacro{\thIn}{\th}
        \pgfmathsetmacro{\thMid}{\th + 0.22}
        \pgfmathsetmacro{\thOut}{\th - 0.04}
        \pgfmathsetmacro{\rMid}{0.42*\rOut + 0.58*\rIn}
        \pgfmathsetmacro{\xIn}{\rIn * cos(\thIn r)}
        \pgfmathsetmacro{\yIn}{-\rIn * sin(\thIn r)}
        \pgfmathsetmacro{\xMid}{\rMid * cos(\thMid r)}
        \pgfmathsetmacro{\yMid}{-\rMid * sin(\thMid r)}
        \pgfmathsetmacro{\xOut}{\rOut * cos(\thOut r)}
        \pgfmathsetmacro{\yOut}{-\rOut * sin(\thOut r)}
        \draw[line width=1.55pt, origcovergreen, line cap=round]
          (\xIn, \yIn) .. controls (\xMid, \yMid) .. (\xOut, \yOut);
      }
    \end{scope}

    % 3. 左侧书名、双细线与出版社（位于螺线大弧上方）
    \draw[line width=0.8pt, origcoverdark] (0.45, 12.66) -- (7.72, 12.66);
    \draw[line width=0.8pt, origcoverdark] (0.45, 10.65) -- (7.38, 10.65);

    \node[anchor=west, inner sep=0pt, text=white] at (0.46, 13.38)
      {\scalebox{1.20}[0.89]{\songti\bfseries\fontsize{25.5}{27}\selectfont 形\hspace{0.23em}形\hspace{0.23em}色\hspace{0.23em}色}};
    \node[anchor=west, inner sep=0pt, text=white] at (0.46, 11.42)
      {\scalebox{1.20}[0.89]{\songti\bfseries\fontsize{25.5}{27}\selectfont 的\hspace{0.23em}曲\hspace{0.23em}线}};
    \node[anchor=west, inner sep=0pt, text=origcoverdark] at (0.50, 9.75)
      {\songti\fontsize{10.0}{11}\selectfont 上\hspace{0.44em}海\hspace{0.44em}教\hspace{0.44em}育\hspace{0.44em}出\hspace{0.44em}版\hspace{0.44em}社};

    % 4. 中下部白色伪球面（对应第 79 节《曳物线与伪球面》：半支曳物线旋转曲面及其上的罗氏几何测地线）
    \begin{scope}
      % 4.1 底部承托弧碗（圆润下探至 y=0.80，精确匹配 600dpi 原封面轮廓）
      \def\bowlPath{
        (4.74, 2.13)
        .. controls (5.16, 1.24) and (6.30, 0.80) .. (8.22, 0.80)
        .. controls (10.14, 0.80) and (11.28, 1.24) .. (11.70, 2.13)
        -- cycle
      }
      \fill[origivory] \bowlPath;
      \begin{scope}
        \clip \bowlPath;
        \shade[left color=origivory, right color=origivoryshade, opacity=0.75]
          (8.0, 0.75) rectangle (11.8, 2.20);
      \end{scope}

      % 4.2 伪球面主体轮廓：精确拟合 600dpi 原封面曳物线半宽与平滑外撇裙底（消除 y=3.0 处微小折点）
      \def\psPath{
        (7.94, 12.74)
        .. controls (7.97, 12.92) and (8.47, 12.92) .. (8.50, 12.74)
        .. controls (8.66, 12.05) and (8.82, 11.05) .. (8.98, 10.00)
        .. controls (9.22, 8.55) and (9.52, 7.15) .. (9.96, 6.00)
        .. controls (10.50, 4.62) and (11.38, 3.60) .. (12.46, 2.94)
        .. controls (12.74, 2.77) and (12.88, 2.58) .. (12.76, 2.38)
        .. controls (12.64, 2.21) and (12.20, 2.15) .. (11.50, 2.06)
        .. controls (10.15, 1.88) and (9.15, 1.81) .. (8.22, 1.81)
        .. controls (7.29, 1.81) and (6.29, 1.88) .. (4.94, 2.06)
        .. controls (4.24, 2.15) and (3.80, 2.21) .. (3.68, 2.38)
        .. controls (3.56, 2.58) and (3.70, 2.77) .. (3.98, 2.94)
        .. controls (5.06, 3.60) and (5.94, 4.62) .. (6.48, 6.00)
        .. controls (6.92, 7.15) and (7.22, 8.55) .. (7.46, 10.00)
        .. controls (7.62, 11.05) and (7.78, 12.05) .. (7.94, 12.74)
        -- cycle
      }

      % 象牙白填色 + 柔和右下部暖象牙微阴影（保持原版纸雕般干净边缘）
      \fill[origivory] \psPath;
      \begin{scope}
        \clip \psPath;
        \shade[left color=origivory, right color=origivoryshade!85!origivory, shading angle=78]
          (6.5, 1.7) rectangle (12.9, 12.9);

        % 4.3 伪球面上的 10 条罗氏几何测地线（逐条对照原封面 600dpi 精确坐标还原）
        \begin{scope}[line width=1.15pt, origcurve, line cap=round]
          % 曲线 1：顶部延伸至中左部的纵贯缓弯测地线
          \draw (8.01, 12.88)
            .. controls (7.99, 11.40) and (7.92, 9.60) .. (7.64, 8.15)
            .. controls (7.36, 6.75) and (6.82, 5.85) .. (6.37, 5.14);

          % 曲线 2：近顶端向右上方微拱的短测地线
          \draw (7.68, 11.55) .. controls (7.88, 11.88) and (8.12, 12.14) .. (8.42, 12.31);

          % 曲线 3：上段向右下方斜穿的测地线
          \draw (7.62, 11.18) .. controls (7.84, 10.82) and (8.12, 10.50) .. (8.42, 10.22);

          % 曲线 4：中上段微倾横向测地线
          \draw (7.24, 8.76) .. controls (7.58, 8.73) and (7.95, 8.65) .. (8.25, 8.51);

          % 曲线 5：中段右倾长斜穿测地线
          \draw (7.08, 7.12) .. controls (7.65, 6.64) and (8.45, 6.12) .. (8.96, 5.53);

          % 曲线 6：中右起向左下大幅扫掠的长测地线
          \draw (8.86, 7.56)
            .. controls (8.62, 6.50) and (8.28, 5.25) .. (7.87, 4.00)
            .. controls (7.52, 3.12) and (7.10, 2.55) .. (6.72, 2.18);

          % 曲线 7：中下段右倾长斜穿测地线
          \draw (7.14, 5.88) .. controls (7.75, 5.10) and (8.38, 4.46) .. (8.99, 3.98);

          % 曲线 8：中轴附近纵下偏左弯的测地线
          \draw (8.30, 5.93)
            .. controls (8.18, 4.80) and (8.16, 3.60) .. (8.07, 3.00)
            .. controls (7.98, 2.52) and (7.82, 2.22) .. (7.65, 2.02);

          % 曲线 9：左下裙部右倾斜穿测地线
          \draw (5.97, 4.42) .. controls (6.52, 3.74) and (7.28, 2.94) .. (8.05, 2.27);

          % 曲线 10：最左侧裙摆 C 形回弯测地线
          \draw (5.11, 3.62) .. controls (4.86, 3.28) and (4.82, 2.82) .. (4.97, 2.44);
        \end{scope}
      \end{scope}

      % 4.4 仅在伪球面底裙前缘与承托弧碗接缝处绘制清晰绿线（主体边缘保持原版干净留白）
      \draw[line width=1.25pt, origcurve!88!origcovergreen, line cap=round]
        (4.18, 2.21)
        .. controls (5.40, 1.98) and (6.80, 1.81) .. (8.22, 1.81)
        .. controls (9.64, 1.81) and (11.04, 1.98) .. (12.26, 2.21);
    \end{scope}
  \end{scope}
\end{tikzpicture}
\clearpage
